\documentclass{article}
\usepackage{geometry}
\usepackage{mathtools}
\usepackage{amsfonts}
\usepackage{amsthm}
\usepackage{amsmath}
\usepackage{todonotes}
\usepackage{enumerate}
\usepackage{graphicx}
\usepackage{subcaption}
\usepackage{url}

\renewcommand{\Re}{\operatorname{Re}}
\renewcommand{\Im}{\operatorname{Im}}
\newcommand{\Log}{\operatorname{Log}}
\newcommand{\Arg}{\operatorname{Arg}}
\newcommand{\C}{\mathbb{C}}
\newcommand{\R}{\mathbb{R}}
\newcommand{\N}{\mathbb{N}}
\newcommand{\eps}{\varepsilon}
\newcommand{\Li}{\operatorname{Li}}

\newcommand\numberthis{\addtocounter{equation}{1}\tag{\theequation}}

\newtheorem{theorem}{Theorem}
\newtheorem{lemma}{Lemma}


\begin{document}
\title{Partial summation and the prime number theorem}
\author{Alexander Dobner}
\date{}
\maketitle

Consider the problem of estimating a sum over primes: $\sum_{a < p \leq b} f(p)$ where $f$ is a continuous function.

By Riemann-Stieltjes integration (or summation by parts), we see
\begin{align*}
\sum_{a < p \leq b} f(p) &= \int_a^b f(t) \, d\pi(t) \\
&= \int_a^b f(t) \, d\Li(t) + \int_a^b f(t) \, d(\pi(t)-\Li(t)) \\
&= \int_a^b \frac{f(t)}{\log t} \, dt + O\left(R(b) |f(b)| + R(a)|f(a)| +\int_a^b R(t) |f'(t)|\, dt\right)
\end{align*}
where $\pi(x) = \sum_{p\leq x} 1$ and $R(x) = | \pi(x)-\Li(x)|$. This estimate is useful if $f$ is rather slowly varying, so that the integral in the error term is small.

Various forms of the prime number theorem give bounds on $R(x)$, the most common of which are
\begin{equation*}
R(x) \ll
\begin{cases}
x \exp(-c \sqrt{\log x})  &\text{Classical form of PNT} \\
\frac{x}{\log^A x} &\text{Weaker (but simpler) PNT} \\
x \exp\left(-c (\log x)^{3/5} (\log\log x)^{-1/5}\right) &\text{Vinogradov-Korobov PNT} \\
\sqrt{x} \log x &\text{Riemann Hypothesis}.
\end{cases}
\end{equation*}

\subsection*{Example}
Suppose $f(t) = \frac{1}{t \log t}$, then $f'(t) \ll \frac{1}{t^2 \log t}$. Applying
\begin{align*}
\sum_{a < p \leq b} \frac{1}{p \log p} &= \frac{1}{\log a} - \frac{1}{\log b} + O_A\left(\log^{-A }a\right),
\end{align*}
which implies
\[
\sum_{p \leq x}\frac{1}{p \log p} = c+\frac{1}{\log x} + O_A\left(\log^{-A} x\right)
\]
\end{document}
